\chapter{Abdominal Mass}

\section*{Definition}
An abdominal mass is an abnormal lump, enlargement, or swelling in the abdomen or abdominal wall, detected on physical examination or imaging. It may arise from an organ, blood vessel, bowel, abdominal wall, or pelvic structure. The clinical significance depends on its location, size, consistency, mobility, tenderness, pulsatility, associated symptoms, and the person\'s medical history; malignancy is one possibility but is not the only cause.

\section*{When to See a Doctor}

\begin{itemize}
 \item Any newly discovered or persistent abdominal mass should be assessed by a clinician; concern is higher with weight loss, loss of appetite, anemia, or increasing age, but a mass in a younger person also requires evaluation.
 \item A rapidly growing mass, increasing abdominal girth, severe or worsening pain, or a mass that is hard, fixed, or very tender
 \item Vomiting, inability to pass stool or gas, GI bleeding, jaundice, fainting, or difficulty breathing accompanying the mass
 \item A pulsatile mass, especially with abdominal or back pain, collapse, sweating, or weakness; call emergency services and do not drive yourself
 \item Fever or feeling very unwell, night sweats, or unexplained weight loss accompanying the mass
\end{itemize}

\section*{How It Develops}
Abdominal masses arise from organomegaly (liver, spleen, kidney), neoplastic growth (primary or metastatic), fluid collections (abscess, cyst), or anatomic abnormalities (aneurysm, hernia). The mass effect may compress adjacent structures, impair organ function, or cause obstructive symptoms.

\section*{Causes}
\begin{itemize}
 \item \textbf{Right upper quadrant:} Hepatomegaly (cirrhosis, hepatitis, malignancy), gallbladder hydrops or empyema, hepatic adenoma, hepatocellular carcinoma, colonic mass
 \item \textbf{Epigastric:} Pancreatic pseudocyst, pancreatic mass, gastric tumor, abdominal aortic aneurysm, left-lobe hepatomegaly, or a prominent aortic pulsation that is not an aneurysm
 \item \textbf{Left upper quadrant:} Splenomegaly (portal hypertension, lymphoma, infection, hemolytic anemia), polycystic kidney, gastric or colonic mass
 \item \textbf{Flank/Lumbar:} Renal mass (hypernephroma, hydronephrosis, polycystic kidney), adrenal mass (pheochromocytoma, adenoma), psoas abscess
 \item \textbf{Pelvic:} Ovarian mass (cyst, neoplasm), uterine fibroids, pelvic kidney, bladder distension, diverticular abscess
 \item \textbf{Diffuse/generalized:} Carcinomatosis, tuberculosis, lymphoma, intra-abdominal abscess, massive ascites
\end{itemize}

\section*{Related Symptoms}
\begin{itemize}
 \item Abdominal pain, distension, early satiety, nausea, vomiting, constipation, weight loss, palpable lump, jaundice
\end{itemize}
\textbf{Important:} A pulsatile or expansile abdominal mass with abdominal or back pain may indicate an abdominal aortic aneurysm or another vascular emergency. Seek emergency care immediately. Do not repeatedly palpate the mass or wait for a routine appointment.

\section*{Medical Evaluation}
\begin{itemize}
 \item Comprehensive history and physical examination determine whether the mass is in the abdominal wall or deeper, and assess its location, size, consistency, mobility, tenderness, pulsatility, relation to respiration, percussion findings, hernias, lymph nodes, and signs of obstruction or systemic illness.
 \item Ultrasound is often an appropriate initial test to distinguish solid from cystic tissue, assess size and vascularity, identify an organ of origin, and evaluate pelvic or superficial masses. It is also used urgently when a vascular cause is suspected.
 \item CT of the abdomen and pelvis, usually with an appropriate contrast protocol, may be needed for characterization, assessment of local spread or obstruction, staging, and treatment planning. The choice of contrast considers kidney function, allergy history, and pregnancy.
 \item MRI is useful for problem-solving in selected liver, pancreatic, pelvic, abdominal-wall, and soft-tissue lesions, and when avoiding ionizing radiation is important.
 \item Laboratory tests are guided by the suspected cause and may include a complete blood count, kidney and liver tests, electrolytes, inflammatory markers, urinalysis, pregnancy testing when relevant, and disease-specific tests. Tumor markers such as CEA, CA 19-9, AFP, or CA-125 are not general screening tests and should be ordered selectively; they do not by themselves diagnose or exclude cancer.
 \item Image-guided biopsy may be needed for tissue diagnosis when malignancy or another condition requiring histology is suspected. Biopsy should be planned by the treating team after imaging because some masses are diagnosed or treated surgically without biopsy, and an improperly placed biopsy can complicate later surgery.
\end{itemize}

\section*{Treatment Options}
\begin{itemize}
 \item Treatment is directed by cause and may include surgery, systemic therapy, radiation, or active surveillance for selected benign or malignant masses. Cancer treatment depends on the tumor type, stage, symptoms, and overall health.
 \item Observation with planned follow-up and serial imaging may be appropriate for selected benign, asymptomatic lesions such as a simple renal cyst or a typical hepatic hemangioma.
 \item Abscesses may require antibiotics and drainage, either percutaneously or surgically, depending on their size, location, and clinical severity.
 \item A suspected abdominal aortic aneurysm requires urgent vascular assessment. Repair decisions depend on symptoms, rupture, aneurysm size and growth, anatomy, and individual risk; a symptomatic or ruptured aneurysm is an emergency and should not be managed by a size threshold at home.
 \item Palliative measures such as drainage, stenting, debulking in selected cases, and pain control may be used for advanced or unresectable disease.
\end{itemize}

\section*{Self-Care and Home Management}
\begin{itemize}
 \item Avoid aggressive palpation or manipulation of a known abdominal mass
 \item Keep a symptom diary noting changes in size, pain, bowel habits, and weight
 \item Eat and drink as tolerated unless vomiting, obstruction, or a clinician\'s instructions require otherwise; a bland diet is not a treatment for an abdominal mass.
 \item If a vascular emergency is suspected, seek emergency care rather than attempting home management or strenuous activity while waiting for help.
 \item Seek medical evaluation for any new or changing abdominal mass rather than observing at home
\end{itemize}

\section*{References}
\begin{thebibliography}{99}
\bibitem{abdominal_mass:ref1} American College of Radiology. ACR Appropriateness Criteria: Palpable Abdominal Mass--Suspected Neoplasm. Revised 2019.
\bibitem{abdominal_mass:ref2} American College of Radiology. ACR Appropriateness Criteria: Pulsatile Abdominal Mass--Suspected Abdominal Aortic Aneurysm. Current online guidance.
\bibitem{abdominal_mass:ref3} National Institute for Health and Care Excellence. Suspected Cancer: Recognition and Referral. NICE guideline NG12, updated 2025.
\bibitem{abdominal_mass:ref4} Society for Vascular Surgery. Clinical Practice Guidelines on the Care of Patients with an Abdominal Aortic Aneurysm. \textit{J Vasc Surg}. 2018;67(1):2--77.e2.
\bibitem{abdominal_mass:ref5} Feldman M, Friedman LS, Brandt LJ, eds. \textit{Sleisenger and Fordtran\'s Gastrointestinal and Liver Disease}, 11th ed. Elsevier, 2021.
\end{thebibliography}
